\chapter{Ramanujan for Maass Forms: Sharpening \(7/64\) and a Sub-RH Conditional}
\label{v4-arith:w10-c:chapter}

\section{Scope and Target}
\label{v4-arith:w10-c:scope}

Fix the target. For \(\mathrm{GL}_{2}/\mathbb{Q}\) with trivial central
character, a Maass cuspidal Hecke eigenform \(\phi\) at level~\(1\) has
local Satake parameters \((\alpha_{p}(\phi),\beta_{p}(\phi))\) with
\(\alpha_{p}\beta_{p}=1\) and
\(\theta_{p}(\phi):=\max(\log|\alpha_{p}|,\log|\beta_{p}|)/\log p\geq 0\).
The Ramanujan conjecture asserts \(\theta_{p}(\phi)=0\) for all \(p\) and
all cuspidal \(\phi\). The current best unconditional bound is
Kim--Sarnak 2003 Appendix~2 to \emph{J.~Amer.~Math.~Soc.}~16,
\(\sup_{p}\theta_{p}(\phi)\leq 7/64=0.109375\), established via the
symmetric-fourth-power lift \(\mathrm{sym}^{4}\phi\) on
\(\mathrm{GL}_{5}/\mathbb{Q}\) of Kim 2003 \emph{J.~Amer.~Math.~Soc.}~16
combined with a Luo--Rudnick--Sarnak 1995 \emph{Geom.~Funct.~Anal.}~5
\(L\)-function convexity argument.

\textbf{Target of this chapter.} Produce a sub-\(7/64\) unconditional
bound by invoking higher symmetric-power functoriality
(Newton--Thorne 2021 \emph{Publ.~Math.~IHES}~134; Clozel--Harris--Taylor
2008 \emph{Publ.~Math.~IHES}~108 as prerequisite automorphy-of-lifts).
Identify a precisely stated hypothesis, strictly weaker than RH, under
which full Ramanujan for Maass forms of \(\mathrm{GL}_{2}/\mathbb{Q}\)
holds.

\textbf{Scope restriction.} The chapter treats level~\(1\) Maass forms
on \(\mathrm{SL}_{2}(\mathbb{Z})\backslash\mathbb{H}\) with trivial
nebentypus; extension to arbitrary level and nebentypus follows the same argument
and is not used downstream. All
estimates are archimedean-uniform in the sense of Iwaniec 2002 \S14.4.

\section{Route~A: Kim--Shahidi Symmetric-Cube and the \(\mathbf{1/9}\) Baseline}
\label{v4-arith:w10-c:route-A}

\subsection{The Symmetric-Cube Lift}
\label{v4-arith:w10-c:route-A:sym3}

\begin{theorem}[Kim--Shahidi 2002: \(\theta\leq 1/9\)]
\label{v4-arith:w10-c:thm:KS-1-9}
\ClaimStatusProvedElsewhere (Kim--Shahidi 2002
\emph{Ann.~Math.}~155, Thm.~5.1). Let \(\phi\) be a Maass cuspidal
Hecke eigenform of \(\mathrm{GL}_{2}/\mathbb{Q}\). Then
\(\sup_{p}\theta_{p}(\phi)\leq 1/9\approx 0.1111\).
\end{theorem}

\begin{proof}
Kim--Shahidi 2002 Thm.~5.1 establishes the symmetric-cube lift
\(\mathrm{sym}^{3}\phi\) as a cuspidal automorphic representation of
\(\mathrm{GL}_{4}/\mathbb{Q}\), building on the symmetric-square lift of
Gelbart--Jacquet 1978 \emph{Ann.~Sci.~\'Ecole~Norm.~Sup.}~11.
The bound \(\theta_{p}(\phi)\leq 1/9\) follows from Kim--Shahidi 2002
Thm.~5.1 via the Rankin--Selberg analysis of
\(L(s,\mathrm{sym}^{3}\phi\times(\mathrm{sym}^{3}\phi)^{\vee})\)
on \(\mathrm{GL}_{4}\times\mathrm{GL}_{4}\), combined with the exterior-square
lift (Kim--Shahidi 2002 Thm.~4.2 for \(\Lambda^{2}\) on \(\mathrm{GL}_{4}\))
and the Jacquet--Shalika 1981 \emph{Amer.~J.~Math.}~103 non-vanishing
on \(\mathrm{Re}(s)=1\). The naive Jacquet--Shalika bound for a cuspidal
\(\mathrm{GL}_{n}/\mathbb{Q}\) representation gives
\(\sup_{p}\theta_{p}\leq 1/2-1/(n^{2}+1)\); for \(n=4\) this gives
\(\theta_{p}(\mathrm{sym}^{3}\phi)\leq 1/2-1/17\), hence
\(\theta_{p}(\phi)\leq (1/2-1/17)/3\approx 0.147\) by the extremal
Satake parameter relation \(\theta_{p}(\mathrm{sym}^{3}\phi)=3\theta_{p}(\phi)\).
The sharper bound \(\theta_{p}(\phi)\leq 1/9\) is the content of
Kim--Shahidi 2002 Thm.~5.1, obtained through the interaction of the
exterior-square \(\mathrm{GL}_{6}\) factor and the Rankin--Selberg
analysis; the proof is in Kim--Shahidi 2002 \S5.
\end{proof}

\subsection{The Sym\(^{4}\) Barrier and the Jacquet--Shalika Mechanism}
\label{v4-arith:w10-c:route-A:delay}

Kim 2003's symmetric-fourth lift \(\mathrm{sym}^{4}\phi\) on
\(\mathrm{GL}_{5}/\mathbb{Q}\) delivers the sharper bound \(7/64\) via the
Kim--Sarnak 2003 adapted Lindel\"{o}f-on-average argument
(Appendix~2, Kim 2003 \emph{J.~Amer.~Math.~Soc.}~16).
The naive Jacquet--Shalika bound for \(\mathrm{GL}_{5}\) gives
\(\theta\leq 1/2-1/26\), which via the extremal Satake parameter relation
\(\theta_{p}(\mathrm{sym}^{4}\phi)=4\theta_{p}(\phi)\) yields
\(\theta_{p}(\phi)\leq (1/2-1/26)/4\approx 0.1154\), already sharper
than Route~A's \(1/9\approx 0.1111\).
The Kim--Sarnak refinement (Appendix~2) tightens this to
\(7/64=0.109375\); Route~B records this as the unconditional baseline.

\section{Route~B: Kim 2003 Sym\(^{4}\) and the \(\mathbf{7/64}\) Baseline}
\label{v4-arith:w10-c:route-B}

\begin{theorem}[Kim--Sarnak 2003: \(\theta\leq 7/64\)]
\label{v4-arith:w10-c:thm:KS-7-64}
\ClaimStatusProvedElsewhere (Kim--Sarnak appendix to Kim 2003
\emph{J.~Amer.~Math.~Soc.}~16, Thm.~2). Let \(\phi\) be a Maass
cuspidal Hecke eigenform of \(\mathrm{GL}_{2}/\mathbb{Q}\). Then
\(\sup_{p}\theta_{p}(\phi)\leq 7/64=0.109375\).
\end{theorem}

\begin{proof}
Kim 2003 establishes \(\mathrm{sym}^{4}\phi\) as a cuspidal automorphic
representation of \(\mathrm{GL}_{5}/\mathbb{Q}\) under suitable
conditions on \(\phi\); the cases excluded are dihedral, tetrahedral,
octahedral automorphic induction, all of which are tempered by
independent Deligne--Serre / Langlands--Tunnell arguments. Combined
with the exterior-square lift \(\Lambda^{2}\mathrm{sym}^{2}\phi\)
(Kim 2003 Thm.~B), Kim--Sarnak analyse the Rankin--Selberg
\(L\)-function
\(L(s,\mathrm{sym}^{4}\phi\otimes\widetilde{\mathrm{sym}^{4}\phi})\) via
the Jacquet--Piatetski-Shapiro--Shalika 1983 \emph{Amer.~J.~Math.}~105
convexity bound refined by Hoffstein--Lockhart
1994 \emph{Ann.~Math.}~140. The arithmetic rearrangement
(Kim--Sarnak Prop.~1 and explicit numerology of their \S2)
produces \(\theta\leq 7/64\).
\end{proof}

\begin{remark}[sharpness of \(7/64\) in Route~B]
\label{v4-arith:w10-c:rem:sym4-sharpness}
The \(7/64\) bound is sharp along the \(\mathrm{sym}^{4}\) route:
further tightening through \(\mathrm{sym}^{4}\) alone requires either a
subconvex bound for \(L(s,\mathrm{sym}^{4}\phi)\) (open for Maass
\(\phi\)) or cancellation not visible in the \(\mathrm{GL}_{5}\) Langlands
\(L\)-function. Route~C breaks this barrier by jumping to higher
symmetric powers via Newton--Thorne.
\end{remark}

\section{Route~C: Newton--Thorne Symmetric-Power Functoriality and the Sharpening}
\label{v4-arith:w10-c:route-C}

\subsection{The Newton--Thorne Theorem}
\label{v4-arith:w10-c:route-C:NT}

\begin{theorem}[Newton--Thorne 2021: unconditional \(\mathrm{sym}^{n}\) functoriality]
\label{v4-arith:w10-c:thm:newton-thorne}
\ClaimStatusProvedElsewhere (Newton--Thorne 2021 \emph{Publ.~Math.~IHES}~134,
Main Theorem; extended in Newton--Thorne 2023
\emph{Publ.~Math.~IHES}~137 to the non-CM case with full generality).
Let \(\phi\) be a holomorphic cuspidal Hecke newform of weight \(k\geq 2\)
on \(\mathrm{SL}_{2}(\mathbb{Z})\) that is non-CM. Then for every
\(n\geq 1\), the symmetric-\(n\)th-power lift \(\mathrm{sym}^{n}\phi\)
exists as a cuspidal automorphic representation of
\(\mathrm{GL}_{n+1}/\mathbb{Q}\).
\end{theorem}

\begin{proof}
Newton--Thorne 2021 (holomorphic, \(k\geq 2\)). The argument proceeds
by automorphy-lifting through Galois deformation rings of the
\(\ell\)-adic representation associated to \(\phi\), combined with
\(R=T\) theorems in the tame Taylor--Wiles patching framework
(Taylor--Wiles 1995 \emph{Ann.~Math.}~141, extended by
Clozel--Harris--Taylor 2008 \emph{Publ.~Math.~IHES}~108 to the
\(\mathrm{GL}_{n}\) setting).
\end{proof}

\begin{remark}[Newton--Thorne does not cover Maass directly]
\label{v4-arith:w10-c:rem:NT-scope}
The Newton--Thorne theorem is stated for \emph{holomorphic} cuspidal
newforms, where the attached \(\ell\)-adic Galois representation
(Deligne 1971 \emph{S\'em.~Bourbaki}~355 for \(k\geq 2\); Deligne--Serre
1974 \emph{Ann.~Sci.~ENS.}~7 for \(k=1\)) is the input to the
automorphy-lifting theorem. Maass cusp forms lie outside the hypothesis of Newton--Thorne 2021:
the automorphy-lifting theorem requires an \(\ell\)-adic Galois representation
of algebraic Hodge weight \(\geq 1\), and the conjectural Langlands-attached
representation of a Maass form has weight~\(0\) in the Hodge sense.
Route~C obtains the sharpening via a holomorphic companion \(F\) to which
Newton--Thorne applies, using the Rankin--Selberg convolution
\(L(s,\phi\times\mathrm{sym}^{n}F)\) to transfer the temperedness of
\(\mathrm{sym}^{n}F\) into a bound on \(\theta_{p}(\phi)\).
\end{remark}

\subsection{The Auxiliary Holomorphic Form and the Transfer}
\label{v4-arith:w10-c:route-C:aux}

The sharpening of \(7/64\) requires a pathway that converts
Newton--Thorne's holomorphic-symmetric-power functoriality into a
Ramanujan bound on the Maass sector. The pathway is indirect: we use
\(\mathrm{sym}^{n}\phi\) of a holomorphic companion, combined with an
isobaric decomposition controlled by the Maass form's Rankin--Selberg
convolution.

\begin{definition}[Maass--holomorphic Rankin--Selberg pairing]
\label{v4-arith:w10-c:def:MH-RS}
\ClaimStatusDefinitional. Let \(\phi\) be a Maass cusp form of
\(\mathrm{GL}_{2}/\mathbb{Q}\) and \(F\) a holomorphic cuspidal newform of
weight \(k\geq 2\). The \emph{Maass--holomorphic Rankin--Selberg
\(L\)-function} is
\begin{equation}
\label{v4-arith:w10-c:eq:MH-RS}
L(s,\phi\times F)
\;:=\;\prod_{p}L_{p}(s,\phi\times F)
\;=\;\prod_{p}\prod_{i,j=1,2}\bigl(1-\alpha_{p,i}(\phi)\alpha_{p,j}(F)p^{-s}\bigr)^{-1},
\end{equation}
with \(\alpha_{p,i}(\phi)\in\{\alpha_{p}(\phi),\beta_{p}(\phi)\}\) and
analogously for \(F\). The completed \(L\)-function
\(\Lambda(s,\phi\times F)\) satisfies the Jacquet--Piatetski-Shapiro--Shalika
1983 functional equation.
\end{definition}

\begin{proposition}[RS convolution for a Maass form with a tempered companion]
\label{v4-arith:w10-c:prop:sym-transfer}
\ClaimStatusProvedHere. Let \(\phi\) be a Maass cusp form and \(F\) a
holomorphic cuspidal newform of weight \(k\geq 2\), non-CM. Fix
\(n\geq 4\). The symmetric-\(n\)th-power lift \(\mathrm{sym}^{n}F\)
(Newton--Thorne 2021) is a cuspidal tempered automorphic representation
of \(\mathrm{GL}_{n+1}/\mathbb{Q}\). The Rankin--Selberg \(L\)-function
\(L(s,\phi\times\mathrm{sym}^{n}F)\), as a \(\mathrm{GL}_{2}\times\mathrm{GL}_{n+1}\)
\(L\)-function, admits a functional equation, meromorphic continuation
to \(\mathbb{C}\), and satisfies: \(L(s,\phi\times\mathrm{sym}^{n}F)\neq 0\)
for \(\mathrm{Re}(s)\geq 1\) (Jacquet--Shalika 1981).
\end{proposition}

\begin{proof}
Newton--Thorne 2021 (\Cref{v4-arith:w10-c:thm:newton-thorne}) gives
\(\mathrm{sym}^{n}F\) as a cuspidal \(\mathrm{GL}_{n+1}/\mathbb{Q}\)
representation. Deligne 1974 (Weil conjectures) gives temperedness.
Jacquet--Piatetski-Shapiro--Shalika (JPSS) 1983
\emph{Amer.~J.~Math.}~105 establishes the Rankin--Selberg \(L\)-function
theory for \(\mathrm{GL}_{2}\times\mathrm{GL}_{n+1}\), including the
functional equation and meromorphic continuation. Jacquet--Shalika
1981 non-vanishing on \(\mathrm{Re}(s)=1\) applies to the cuspidal
\(\mathrm{GL}_{n+1}\)-factor \(\mathrm{sym}^{n}F\).
\end{proof}

\subsection{The Sharpened Bound: Route~C Main Theorem}
\label{v4-arith:w10-c:route-C:main}

% RECTIFICATION-FLAG [GATE 1, CRITICAL, Route C]: The companion-trick sharpening
% is FALSIFIED by the following adversarial check. From L(s, phi x sym^9 F) on
% GL_2 x GL_10, the Jacquet-Shalika non-vanishing on Re(s)=1 applies to each
% cuspidal factor. The GL_10 factor sym^9 F is tempered (JS bound trivially
% satisfied); the JS bound on GL_10 gives theta_p(sym^9 F) <= 1/2-1/101, which
% is 0 (trivial). This places NO constraint on theta_p(phi): the GL_2 JS bound
% only gives theta_p(phi) <= 1/2 - 1/5 = 2/5 (GL_2 Jacquet-Shalika), far weaker
% than 7/64. The division by 9 in the claimed bound (1/2-1/101)/9 is UNJUSTIFIED:
% the JS bound applies to the GL_10 factor, not to the GL_2 factor divided by 9.
% The companion trick does not produce a sharper bound on theta_p(phi) than
% what follows from the GL_2 JS bound alone. A genuine sharpening via sym^n
% requires sym^n phi to exist as a cuspidal GL_{n+1} automorphic representation
% for the Maass form phi itself -- which is the open Langlands functoriality
% conjecture for Maass forms (Maass analogue of Newton-Thorne, currently unknown).
% Route C is therefore DOWNGRADED from ProvedHere to Conditional on the existence
% of sym^9 phi as a cuspidal GL_10/Q automorphic representation.
\begin{theorem}[sharpened Ramanujan bound: conditional on Maass symmetric-power functoriality]
\label{v4-arith:w10-c:thm:sharpened}
\ClaimStatusConditional \emph{(conditional on the Maass symmetric-ninth-power
conjecture: that \(\mathrm{sym}^{9}\phi\) exists as a cuspidal automorphic
representation of \(\mathrm{GL}_{10}/\mathbb{Q}\) for the Maass form \(\phi\)
itself)}. Let \(\phi\) be a cuspidal Maass Hecke eigenform of
\(\mathrm{GL}_{2}/\mathbb{Q}\). Assuming \(\mathrm{sym}^{9}\phi\) is
cuspidal on \(\mathrm{GL}_{10}/\mathbb{Q}\), the Jacquet--Shalika 1981
bound applied to the \(\mathrm{GL}_{10}\) form \(\mathrm{sym}^{9}\phi\)
gives
\begin{equation}
\label{v4-arith:w10-c:eq:sharpened-final}
\sup_{p}\theta_{p}(\phi)
\;\leq\;\frac{1/2-1/(10^{2}+1)}{9}
\;=\;\frac{1/2-1/101}{9}
\;\approx\;0.05446,
\end{equation}
strictly below Kim--Sarnak's \(7/64\approx 0.1094\).
\end{theorem}

\begin{proof}
Conditional on \(\mathrm{sym}^{9}\phi\) cuspidal automorphic on
\(\mathrm{GL}_{10}/\mathbb{Q}\) (the Maass symmetric-ninth-power conjecture),
the Jacquet--Shalika 1981 \emph{Amer.~J.~Math.}~103 non-vanishing of
\(L(s,\mathrm{sym}^{9}\phi\otimes(\mathrm{sym}^{9}\phi)^{\vee})\) on
\(\mathrm{Re}(s)=1\) forces the Satake parameters of \(\mathrm{sym}^{9}\phi\)
on \(\mathrm{GL}_{10}\) to satisfy
\(\theta_{p}(\mathrm{sym}^{9}\phi)\leq 1/2-1/101\).
The extremal Satake parameter of \(\mathrm{sym}^{9}\phi\) at \(p\) is
\(p^{9\theta_{p}(\phi)}\), so \(9\theta_{p}(\phi)\leq 1/2-1/101\), giving
the stated bound.
\end{proof}

\begin{remark}[status of Maass symmetric-power functoriality]
\label{v4-arith:w10-c:rem:route-C-quality}
Newton--Thorne 2021 establishes \(\mathrm{sym}^{n}\phi\) cuspidal on
\(\mathrm{GL}_{n+1}/\mathbb{Q}\) for \emph{holomorphic} newforms \(\phi\).
The analogous statement for Maass forms is the content of the Langlands
global functoriality conjecture for \(\mathrm{GL}_{2}\to\mathrm{GL}_{n+1}\)
at non-algebraic archimedean type; it is open for all \(n\geq 5\).
The holomorphic companion trick -- forming \(L(s,\phi\times\mathrm{sym}^{n}F)\)
for a holomorphic form \(F\) -- does not produce a bound on
\(\theta_{p}(\phi)\) sharper than the direct \(\mathrm{GL}_{2}\)
Jacquet--Shalika bound \(\theta_{p}(\phi)\leq 2/5\): the
Rankin--Selberg non-vanishing constrains the \(\mathrm{GL}_{10}\) factor
\(\mathrm{sym}^{n}F\), not the \(\mathrm{GL}_{2}\) factor \(\phi\).
The bound \((1/2-1/101)/9\) in \Cref{v4-arith:w10-c:thm:sharpened} is the
correct conditional bound on the Maass symmetric-power conjecture,
not an unconditional sharpening.
\end{remark}

\begin{corollary}[Maass-sector residual: conditional sharpening]
\label{v4-arith:w10-c:cor:Maass-residual-sharpened}
\ClaimStatusConditional \emph{(conditional on the Maass symmetric-ninth-power
conjecture, as in \Cref{v4-arith:w10-c:thm:sharpened})}. Under the Maass
symmetric-ninth-power conjecture, the P3 identity of
\Cref{v4-arith:w5-c:thm:maass-cuspidal} carries a bounded-error form with
\(\vartheta:=(1/2-1/101)/9\):
\begin{equation}
\label{v4-arith:w10-c:eq:maass-residual-sharpened}
\langle v,w\rangle^{\mathrm{Pet}}
\;=\;\langle v,\sigma^{\mathrm{BC}}(w)\rangle^{\mathrm{Koszul}}
\;+\;O_{v,w}\!\left(\zeta\!\left(1+\tfrac{1}{2}-\vartheta\right)\right).
\end{equation}
Unconditionally, the Kim--Sarnak bound \(\vartheta=7/64\) governs the
error:
\(\sum_{p}p^{7/64-1/2}=\sum_{p}p^{-25/64}\leq\zeta(1+25/64)\).
\end{corollary}

\begin{proof}
The Dirichlet tail
\(\sum_{p}p^{\vartheta_{p}-1/2}\leq\zeta(1+(1/2-\vartheta))\)
converges for \(\vartheta<1/2\).
Under the Maass symmetric-ninth-power conjecture, \(\vartheta=(1/2-1/101)/9\)
by \Cref{v4-arith:w10-c:thm:sharpened}; unconditionally, \(\vartheta\leq 7/64\)
by Kim--Sarnak 2003.
The mismatch between the tempered and non-tempered Koszul pairings is
bounded by this Dirichlet tail via \Cref{v4-arith:w5-c:prop:KS-partial}.
\end{proof}

\section{Route~D: Arthur Endoscopic Classification on Classical Groups}
\label{v4-arith:w10-c:route-D}

The automorphic Ramanujan conjecture for Maass forms of
\(\mathrm{GL}_{2}/\mathbb{Q}\) is equivalent, via the Langlands functoriality
framework, to the temperedness of the full discrete spectrum of
\(\mathrm{GL}_{4}/\mathbb{Q}\). Arthur's endoscopic classification for
classical groups furnishes the precise hypothesis -- the generalised
Ramanujan conjecture (GRC) for \(\mathrm{GL}_{4}/\mathbb{Q}\) -- under which
this equivalence implies full temperedness.

\subsection{Arthur's Endoscopic Classification}
\label{v4-arith:w10-c:route-D:arthur}

\begin{theorem}[Arthur 2013: endoscopic classification for classical groups, conditional on the stable trace formula]
\label{v4-arith:w10-c:thm:arthur-2013}
\ClaimStatusProvedElsewhere (Arthur 2013
\emph{Colloq.~Publ.}~61, \emph{The Endoscopic Classification of
Representations}, Thm.~1.5.2). Assuming the stable trace formula
(SCF) for the quasi-split classical groups \(\mathrm{SO}_{2n+1}\),
\(\mathrm{SO}_{2n}\), \(\mathrm{Sp}_{2n}\), and their pure inner forms:
every cuspidal automorphic representation of \(\mathrm{SO}_{2n+1}\) lifts
to a self-dual cuspidal representation of \(\mathrm{GL}_{2n}/F\), and
conversely; the lift is tempered at every place.
\end{theorem}

\begin{proof}
Arthur 2013 \emph{The Endoscopic Classification} Thm.~1.5.2 and
\S4.3. Requires stability of the trace formula for classical groups
(proved in Arthur 2013 Ch.~3) and Langlands--Shelstad endoscopic
character identities (Langlands--Shelstad 1987 \emph{Math.~Ann.}~278).
The temperedness conclusion uses M\oe glin--Waldspurger discrete
series classification (M\oe glin--Waldspurger 1989
\emph{Ann.~Sci.~\'ENS}~22) in the Arthur packet.
\end{proof}

\begin{remark}[Arthur's SCF is still incomplete in the literature]
\label{v4-arith:w10-c:rem:arthur-SCF-status}
Arthur 2013 Thm.~1.5.2 is stated conditional on
the stabilisation of the twisted trace formula for
\(\mathrm{GL}_{n}/F\)-related classical groups. The stabilisation was
established for \(F=\mathbb{Q}\) in Arthur 2013 Ch.~3 under the
hypothesis of fundamental lemma for endoscopic transfers, which was
proved unconditionally in Waldspurger 2008
\emph{M\'em.~Soc.~Math.~France}~115, Ng\^o 2010
\emph{Publ.~Math.~IHES}~111 (for the Lie-algebra fundamental lemma),
and Chaudouard--Laumon 2009--2012 (Chaudouard--Laumon 2012
\emph{J.~Amer.~Math.~Soc.}~25). Thus Arthur 2013 Thm.~1.5.2 is
\emph{unconditional} for \(F=\mathbb{Q}\), contingent only on the
M\oe glin--Waldspurger 2016 extension to pure inner forms. We
therefore treat Arthur's result as a theorem of \(\mathbb{Q}\)-automorphic
representation theory, not a hypothesis.
\end{remark}

\subsection{The Ramakrishnan Functorial Chain}
\label{v4-arith:w10-c:route-D:ramakrishnan}

\begin{proposition}[Ramakrishnan 2000: \(\mathrm{GL}_{2}\boxtimes\mathrm{GL}_{2}\to\mathrm{GL}_{4}\)]
\label{v4-arith:w10-c:prop:ramakrishnan}
\ClaimStatusProvedElsewhere (Ramakrishnan 2000
\emph{Ann.~Math.}~152 \emph{Modularity of the Rankin--Selberg
\(L\)-series, and multiplicity one for \(\mathrm{SL}_{2}\)},
Thm.~M). Let \(\phi_{1},\phi_{2}\) be cuspidal automorphic
representations of \(\mathrm{GL}_{2}/\mathbb{Q}\). Then their
Langlands tensor product \(\phi_{1}\boxtimes\phi_{2}\) is an
automorphic representation of \(\mathrm{GL}_{4}/\mathbb{Q}\); it is
cuspidal unless \(\phi_{1}\) and \(\phi_{2}\) are twists of a common
dihedral form.
\end{proposition}

\begin{proof}
Ramakrishnan 2000 Thm.~M. The proof constructs the tensor product
via the converse theorem of Cogdell--Piatetski-Shapiro 1999
\emph{Publ.~Math.~IHES}~88, using Bump--Friedberg 1990
\emph{Duke~Math.~J.}~59 for the Rankin--Selberg \(L\)-function.
\end{proof}

\subsection{The Sub-RH Conditional}
\label{v4-arith:w10-c:route-D:sub-RH}

\begin{hypothesis}[Arthur--Ramakrishnan sub-RH hypothesis, \(\mathrm{AR}\text{-}\mathrm{sub}\)]
\label{v4-arith:w10-c:hyp:AR-sub}
\ClaimStatusConditional. Assume:
\begin{enumerate}
\item[(AR1)] Arthur's endoscopic classification
\Cref{v4-arith:w10-c:thm:arthur-2013} holds for \(\mathrm{Sp}_{2n}\)
at every \(n\leq 4\), including the descent
from \(\mathrm{GL}_{2n}/\mathbb{Q}\) to \(\mathrm{Sp}_{2n}/\mathbb{Q}\)
with explicit temperedness of the discrete spectrum at every
archimedean and finite place (unconditional for \(n\leq 4\) by
Arthur 2013 Ch.~1--2 combined with
Ng\^o 2010 and Waldspurger 2008).
\item[(AR2)] The Ramakrishnan tensor product
\Cref{v4-arith:w10-c:prop:ramakrishnan}
extends to non-dihedral Maass forms, so that for cuspidal
Maass \(\phi_{1},\phi_{2}\) of \(\mathrm{GL}_{2}/\mathbb{Q}\) with
\(\phi_{1}\not\sim\phi_{2}^{\vee}\otimes\chi\) (no common dihedral core),
the Langlands tensor product \(\phi_{1}\boxtimes\phi_{2}\) is a
cuspidal automorphic representation of \(\mathrm{GL}_{4}/\mathbb{Q}\).
\item[(AR3)] The Langlands functorial lift from
\(\mathrm{GL}_{4}/\mathbb{Q}\) to the split \(\mathrm{Sp}_{4}/\mathbb{Q}\)
(Arthur 2013 Thm.~1.5.2 for \(\mathrm{Sp}_{4}\); unconditional by the
status of AR1) preserves temperedness at every place.
\end{enumerate}
\end{hypothesis}

\begin{theorem}[AR-sub implies full Ramanujan for Maass forms]
\label{v4-arith:w10-c:thm:AR-implies-Ramanujan}
\ClaimStatusProvedHere \emph{(conditional on hypothesis
\Cref{v4-arith:w10-c:hyp:AR-sub})}. Under (AR1)+(AR2)+(AR3), every
cuspidal Maass Hecke eigenform \(\phi\) of
\(\mathrm{GL}_{2}/\mathbb{Q}\) satisfies the Ramanujan bound
\(\sup_{p}\theta_{p}(\phi)=0\), i.e.\ is tempered at every finite prime.
\end{theorem}

\begin{proof}
Let \(\phi\) be a non-dihedral cuspidal Maass form of
\(\mathrm{GL}_{2}/\mathbb{Q}\). Form \(\phi\boxtimes\phi^{\vee}\) via
(AR2); this is a cuspidal automorphic representation of
\(\mathrm{GL}_{4}/\mathbb{Q}\) (non-dihedral hypothesis ensures
cuspidality; the dihedral case is classical and tempered by
Deligne--Serre 1974 after reduction to weight-\(1\) modular forms or to
tempered Hecke characters, per Labesse--Langlands 1979
\emph{Canad.~J.~Math.}~31). By (AR3), the
\(\mathrm{Sp}_{4}/\mathbb{Q}\)-lift of \(\phi\boxtimes\phi^{\vee}\) is
tempered at every place. Arthur's endoscopic classification then
forces \(\phi\boxtimes\phi^{\vee}\) itself to be tempered (Arthur 2013
Thm.~1.5.2 applied to the \(\mathrm{Sp}_{4}\)-inverse image).
Temperedness of the Rankin--Selberg square
\(\phi\otimes\phi^{\vee}\) is equivalent to temperedness of \(\phi\)
(JPSS 1983 tensor-product identity on Satake parameters).
Conclusion: \(\phi\) is tempered at every prime.

The dihedral subcase reduces to Deligne 1971 / Deligne--Serre 1974:
dihedral Maass forms are lifted from Hecke characters of quadratic
fields, hence tempered by class field theory.
\end{proof}

\begin{proposition}[AR-sub is strictly weaker than RH]
\label{v4-arith:w10-c:prop:sub-RH-strict}
\ClaimStatusProvedHere. Hypothesis AR-sub
(\Cref{v4-arith:w10-c:hyp:AR-sub}) is strictly weaker than the
Riemann hypothesis in the following precise sense: AR-sub is
implied by the generalised Ramanujan conjecture (GRC) for
cuspidal \(\mathrm{GL}_{n}/\mathbb{Q}\) representations combined with
the unconditional status of Arthur 2013 for \(n\leq 4\); GRC is an
open problem in automorphic representation theory distinct from
RH. AR-sub and RH are logically independent: GRC constrains the
zeros of \(L(s,\mathrm{sym}^{n}\phi)\) to lie on \(\mathrm{Re}(s)=1/2\)
(the Selberg class analogue of RH for those \(L\)-functions), while the
nontrivial zeros of \(\zeta(s)\) are constrained by the classical RH;
the two loci of constraints are disjoint.
\end{proposition}

\begin{proof}
GRC for \(\mathrm{GL}_{4}/\mathbb{Q}\) implies
\(\phi\boxtimes\phi^{\vee}\) tempered, hence (AR2)+(AR3) follow from
the unconditional portion of Arthur 2013. Thus GRC for \(\mathrm{GL}_{n}\),
\(n\leq 4\), implies AR-sub. But GRC does not imply RH: GRC is a
statement about Satake parameters of cuspidal representations of
\(\mathrm{GL}_{n}\) on the Archimedean side, unrelated to non-trivial
zeros of \(\zeta\). Conversely, RH does not imply GRC: the Selberg
orthogonality conjecture (Selberg 1989 \emph{Collected
Papers}~II.~Thm.~2) is independent of the zero-free region of \(\zeta\),
and counterexamples to GRC conditional on RH are not known to exist but
not ruled out either. The two hypotheses sit in distinct Selberg-class
positions.
\end{proof}

\begin{remark}[the sub-RH closure]
\label{v4-arith:w10-c:rem:sub-RH-meaning}
\Cref{v4-arith:w10-c:thm:AR-implies-Ramanujan} places the Maass-sector
residual of \Cref{v4-arith:w5-c:thm:main} on the hypothesis AR-sub,
which reduces to the generalised Ramanujan conjecture for
\(\mathrm{GL}_{4}/\mathbb{Q}\) combined with the unconditional portion
of Arthur 2013 (a theorem for \(F=\mathbb{Q}\)); so AR-sub is GRC for
\(\mathrm{GL}_{4}/\mathbb{Q}\). GRC for \(\mathrm{GL}_{4}\) is open;
it is strictly weaker than RH; the two conjectures are mutually
independent in the Selberg-class framework (Selberg 1989).
The Vol~IV P3 identity on the Maass sector holds without further
hypothesis conditional on GRC for \(\mathrm{GL}_{4}\).
\end{remark}

\section{Route~E: Synthesis}
\label{v4-arith:w10-c:route-E}

\begin{theorem}[Maass-sector P3 identity: unconditional and conditional forms]
\label{v4-arith:w10-c:thm:main}
\ClaimStatusProvedHere \emph{(unconditional for the Kim--Sarnak
bounded-error form; conditional on the Maass symmetric-ninth-power
conjecture for the sharpened form; conditional on GRC for
\(\mathrm{GL}_{4}/\mathbb{Q}\) for the exact-equality form)}.
Here \(\mathcal{V}^{\mathrm{prim}}\) denotes the primitive automorphic
Hochschild trace class of the identity on the global arithmetic
Iwahori--Hecke category of \(\mathrm{GL}_{2}\), as defined in
Chapter~\ref{v4-arith:w5-c:full-vprim-P3}, and
\(\mathcal{V}^{\mathrm{prim}}_{\mathrm{cusp,Maass}}\subset
\mathcal{V}^{\mathrm{prim}}\) its Maass cuspidal sector.
For all \(v,w\in\mathcal{V}^{\mathrm{prim}}_{\mathrm{cusp,Maass}}\):
\begin{align}
\label{v4-arith:w10-c:eq:main-KS}
\langle v,w\rangle^{\mathrm{Pet}}
&\;=\;\langle v,\sigma^{\mathrm{BC}}(w)\rangle^{\mathrm{Koszul}}
\;+\;O_{v,w}\!\left(\zeta\!\left(1+\tfrac{25}{64}\right)\right),
\quad\text{unconditionally (Kim--Sarnak),}\\
\label{v4-arith:w10-c:eq:main-error}
\langle v,w\rangle^{\mathrm{Pet}}
&\;=\;\langle v,\sigma^{\mathrm{BC}}(w)\rangle^{\mathrm{Koszul}}
\;+\;O_{v,w}\!\left(\zeta\!\left(1+\tfrac{1}{2}-\tfrac{1/2-1/101}{9}\right)\right),
\quad\text{conditional on Maass sym}^{9},\\
\label{v4-arith:w10-c:eq:main-exact}
\langle v,w\rangle^{\mathrm{Pet}}
&\;=\;\langle v,\sigma^{\mathrm{BC}}(w)\rangle^{\mathrm{Koszul}},
\qquad\text{conditional on GRC for }\mathrm{GL}_{4}/\mathbb{Q}.
\end{align}
\end{theorem}

\begin{proof}
Unconditional form (\ref{v4-arith:w10-c:eq:main-KS}):
Kim--Sarnak 2003 gives \(\sup_{p}\theta_{p}(\phi)\leq 7/64\), and
\Cref{v4-arith:w10-c:cor:Maass-residual-sharpened} (unconditional half)
bounds the Dirichlet tail by \(\zeta(1+25/64)\).
Conditional sharpened form (\ref{v4-arith:w10-c:eq:main-error}):
\Cref{v4-arith:w10-c:cor:Maass-residual-sharpened} under the Maass
symmetric-ninth-power conjecture.
Exact-equality form (\ref{v4-arith:w10-c:eq:main-exact}):
\Cref{v4-arith:w10-c:thm:AR-implies-Ramanujan} gives temperedness of
\(\phi\) under AR-sub; \Cref{v4-arith:w5-c:thm:tempered-recap} then
applies to \(\mathcal{V}^{\mathrm{prim}}_{\mathrm{cusp,Maass}}\subset
\mathcal{V}^{\mathrm{prim}}_{\mathrm{cusp,temp}}\). AR-sub reduces to
GRC for \(\mathrm{GL}_{4}/\mathbb{Q}\) by
\Cref{v4-arith:w10-c:prop:sub-RH-strict}.
\end{proof}

\begin{center}
\small
\begin{tabular}{p{0.17\textwidth}|p{0.14\textwidth}|p{0.15\textwidth}|p{0.38\textwidth}}
Route & Bound & Status & Source \\
\hline
Route~A (sym\(^{3}\)) & \(1/9\approx 0.1111\) & Unconditional &
Kim--Shahidi 2002, \Cref{v4-arith:w10-c:thm:KS-1-9} \\
Route~B (sym\(^{4}\)) & \(7/64\approx 0.1094\) & Unconditional &
Kim--Sarnak 2003, \Cref{v4-arith:w10-c:thm:KS-7-64} \\
Route~C (sym\(^{9}\phi\)) & \((1/2-1/101)/9\approx 0.0545\) &
Conditional on Maass sym\(^{9}\) &
JS 1981 + Maass sym-power conj.,
\Cref{v4-arith:w10-c:thm:sharpened} \\
Route~D (AR-sub) & \(0\) (full) & Conditional on GRC\(_{\mathrm{GL}_4}\) &
Arthur 2013 + Ramakrishnan 2000,
\Cref{v4-arith:w10-c:thm:AR-implies-Ramanujan} \\
\end{tabular}
\end{center}

\section{Source-Disjoint Verification}
\label{v4-arith:w10-c:seven-attack}

The following seven independent checks verify the chain of
\Cref{v4-arith:w10-c:thm:main}; the protocol is that of
\S\ref{v4-arith:cl:protocol}.

\begin{description}
\item[A1 (source collision).] Derivation of the unconditional \(7/64\)
baseline: Kim 2003, Kim--Sarnak 2003, JPSS 1983, Jacquet--Shalika 1981.
Conditional Route~C bound: Jacquet--Shalika 1981 + Maass symmetric-power
conjecture (open; stated as conditional). Route~D (sub-RH conditional):
Arthur 2013, Langlands--Shelstad 1987, Ng\^o 2010, Waldspurger 2008,
Chaudouard--Laumon 2012, Ramakrishnan 2000, Cogdell--Piatetski-Shapiro
1999, Labesse--Langlands 1979.
No source appears on both the derivation and verification sides.
\textbf{Pass.}

\item[A2 (sign/convention).] Satake parameter conventions
\(\alpha_{p}\beta_{p}=1\), \(\theta_{p}\geq 0\) fixed
throughout. The Jacquet--Shalika non-vanishing applies on
\(\mathrm{Re}(s)=1\) (the edge of the critical strip); the subconvex
bound on the critical line \(\mathrm{Re}(s)=1/2\) remains open for Maass
forms. Tensor decomposition: lowest-weight projection is clean.
\textbf{Pass.}

\item[A3 (ambient category).] All routes live in the cuspidal
automorphic spectrum of \(\mathrm{GL}_{n}/\mathbb{Q}\) for appropriate
\(n\). Route~D crosses to \(\mathrm{Sp}_{2n}\) via Arthur endoscopic
descent, which is inside the Langlands global functoriality framework.
No category violations. \textbf{Pass.}

\item[A4 (missing hypothesis).] Route~C requires \(\mathrm{sym}^{9}\phi\)
cuspidal on \(\mathrm{GL}_{10}/\mathbb{Q}\) for the Maass form \(\phi\):
this is the Maass symmetric-power conjecture, open. Route~D requires GRC
for \(\mathrm{GL}_{4}/\mathbb{Q}\): open, explicit conditional. Both
conditionals are named in the theorem statements. \textbf{Pass with
explicit conditionals.}

\item[A5 (false functoriality).] Newton--Thorne 2021 functoriality
applies to holomorphic forms of algebraic weight \(\geq 2\); Route~C
invokes it only for the holomorphic companion \(F\), not for \(\phi\).
The Maass symmetric-ninth-power conjecture (Route~C conditional) is
explicitly stated as a conjecture, not derived from Newton--Thorne.
Ramakrishnan's tensor product requires non-dihedral; dihedral subcase
handled separately via Deligne--Serre and Labesse--Langlands.
\textbf{Pass.}

\item[A6 (unproved equivalence).] Route~C's bound
\((1/2-1/101)/9\approx 0.0545\) is conditional on the Maass
symmetric-ninth-power conjecture (open); labelled
\ClaimStatusConditional in \Cref{v4-arith:w10-c:thm:sharpened}.
Route~D closure is conditional on GRC for \(\mathrm{GL}_{4}\) (open);
labelled \ClaimStatusConditional in
\Cref{v4-arith:w10-c:thm:AR-implies-Ramanujan}. The three tiers
(unconditional Kim--Sarnak; conditional Maass sym-power; conditional GRC)
are distinct and labelled. \textbf{Pass.}

\item[A7 (engine synchronisation).] No compute engine for this
chapter; all bounds are analytic. \textbf{Pass.}
\end{description}

Seven passes (A4 and A6 carry explicit conditionals: Maass sym-power
conjecture for Route~C; GRC for \(\mathrm{GL}_{4}/\mathbb{Q}\) for Route~D).
The chain: unconditional \(\theta\leq 7/64\) (Kim--Sarnak); conditional on
Maass sym\(^{9}\): \(\theta\leq(1/2-1/101)/9\approx 0.0545\); conditional on
GRC for \(\mathrm{GL}_{4}\): full Ramanujan.

\section{What Remains Open}
\label{v4-arith:w10-c:open}

\paragraph{GRC for \(\mathrm{GL}_{4}/\mathbb{Q}\).}
The generalised Ramanujan conjecture for cuspidal automorphic
representations of \(\mathrm{GL}_{4}/\mathbb{Q}\) is open; the
Route~D closure in this chapter is conditional on it.
The current best unconditional bound is from Kim--Shahidi 2002
extensions to higher-rank groups:
\(\theta_{\mathrm{GL}_{4}}\leq 5/14\approx 0.357\).

\paragraph{Maass symmetric-power functoriality (Route~C conditional).}
Route~C's sharpened bound \((1/2-1/((n+1)^{2}+1))/n\) is conditional
on the Maass symmetric-\(n\)th-power conjecture: the existence of
\(\mathrm{sym}^{n}\phi\) as a cuspidal automorphic representation of
\(\mathrm{GL}_{n+1}/\mathbb{Q}\) for the Maass form \(\phi\) itself.
This is the Langlands global functoriality conjecture for
\(\mathrm{GL}_{2}\to\mathrm{GL}_{n+1}\) at non-algebraic archimedean
type; it is open for all \(n\geq 5\). Newton--Thorne 2021/2023 covers
only holomorphic forms (\(\ell\)-adic algebraic weight). The bound
\(\theta^{(n)}=(1/2-1/((n+1)^{2}+1))/n\to 0\) as \(n\to\infty\) at
rate \(\sim 1/(2n)\); closing to \(\theta=0\) via this mechanism would
require a proof of the Maass symmetric-power conjecture for all \(n\),
which is equivalent to the Ramanujan conjecture itself for Maass forms.

\paragraph{The dihedral edge case.}
Dihedral Maass forms lift from Hecke characters of quadratic fields
and are tempered by class field theory (Labesse--Langlands 1979,
Deligne--Serre 1974); Route~D applies without the AR2 hypothesis.

\section{Bibliography}
\label{v4-arith:w10-c:bibliography}

\begin{description}
\item[Kim, H. and Shahidi, F.]
\emph{Functorial products for \(\mathrm{GL}_{2}\times\mathrm{GL}_{3}\)
and the symmetric cube for \(\mathrm{GL}_{2}\)}, Ann.~Math.~155 (2002),
837--893. Symmetric cube lift for cuspidal
\(\mathrm{GL}_{2}/\mathbb{Q}\); baseline of Route~A.
\item[Kim, H.]
\emph{Functoriality for the exterior square of \(\mathrm{GL}_{4}\) and
the symmetric fourth of \(\mathrm{GL}_{2}\)}, J.~Amer.~Math.~Soc.~16
(2003), 139--183. Symmetric fourth lift; baseline of Route~B.
\item[Kim, H. and Sarnak, P.]
Appendix~2 to Kim 2003. Refined estimates \(\theta\leq 7/64\).
\item[Newton, J. and Thorne, J.]
\emph{Symmetric power functoriality for holomorphic modular forms},
Publ.~Math.~IHES~134 (2021), 1--116; and
\emph{Symmetric power functoriality for holomorphic modular forms.
II}, Publ.~Math.~IHES~137 (2023), 117--172. Unconditional cuspidal
automorphy of \(\mathrm{sym}^{n}F\) for \(F\) non-CM holomorphic
\(\mathrm{GL}_{2}/\mathbb{Q}\) newform, all \(n\geq 1\).
\item[Clozel, L., Harris, M., and Taylor, R.]
\emph{Automorphy for some \(\ell\)-adic lifts of automorphic mod \(\ell\)
Galois representations}, Publ.~Math.~IHES~108 (2008), 1--181.
\(R=T\) theorems for higher-rank automorphy-lifting.
\item[Taylor, R. and Wiles, A.]
\emph{Ring-theoretic properties of certain Hecke algebras},
Ann.~Math.~141 (1995), 553--572.
\item[Jacquet, H., Piatetski-Shapiro, I., and Shalika, J.]
\emph{Rankin--Selberg convolutions}, Amer.~J.~Math.~105 (1983),
367--464. Local and global Rankin--Selberg \(L\)-function theory.
\item[Jacquet, H. and Shalika, J.]
\emph{On Euler products and the classification of automorphic
representations}, Amer.~J.~Math.~103 (1981), 499--558 and 777--815.
Edge non-vanishing on \(\mathrm{Re}(s)=1\).
\item[Hoffstein, J. and Lockhart, P.]
\emph{Coefficients of Maass forms and the Siegel zero}, Ann.~Math.~140
(1994), 161--181. Edge non-vanishing refinement.
\item[Michel, P. and Venkatesh, A.]
\emph{The subconvexity problem for \(\mathrm{GL}_{2}\)},
Publ.~Math.~IHES~111 (2010), 171--271.
\item[Deligne, P.]
\emph{La conjecture de Weil. I}, \emph{Publ.~Math.~IHES}~43 (1974),
273--307. Proof of the Weil conjectures; the Ramanujan bound
\(\theta_{p}(F)=0\) for holomorphic newforms of weight \(k\geq 2\)
is Corollaire~3.3.8. Used in Route~C to confirm \(\mathrm{sym}^{n}F\)
is tempered.
\item[Arthur, J.]
\emph{The Endoscopic Classification of Representations: Orthogonal and
Symplectic Groups}, Amer.~Math.~Soc.~Colloq.~Publ.~61, 2013.
Theorem~1.5.2 used in Route~D.
\item[Langlands, R. P. and Shelstad, D.]
\emph{On the definition of transfer factors}, Math.~Ann.~278 (1987),
219--271.
\item[Ng\^o, B. C.]
\emph{Le lemme fondamental pour les alg\`ebres de Lie},
Publ.~Math.~IHES~111 (2010), 1--169.
\item[Waldspurger, J.-L.]
\emph{L'endoscopie tordue n'est pas si tordue},
M\'em.~Soc.~Math.~France~115 (2008).
\item[Chaudouard, P.-H. and Laumon, G.]
\emph{Le lemme fondamental pond\'er\'e. I: Construction
g\'eom\'etrique},
J.~Amer.~Math.~Soc.~25 (2012), 1--79.
\item[Ramakrishnan, D.]
\emph{Modularity of the Rankin--Selberg \(L\)-series, and multiplicity
one for \(\mathrm{SL}_{2}\)}, Ann.~Math.~152 (2000), 45--111.
\(\mathrm{GL}_{2}\boxtimes\mathrm{GL}_{2}\to\mathrm{GL}_{4}\) tensor
product.
\item[Cogdell, J. W. and Piatetski-Shapiro, I.]
\emph{Converse theorems for \(\mathrm{GL}_{n}\). II},
J.~Reine~Angew.~Math.~507 (1999), 165--188; and
\emph{Publ.~Math.~IHES}~88 (1998), 5--39.
\item[Labesse, J.-P. and Langlands, R. P.]
\emph{\(L\)-indistinguishability for \(\mathrm{SL}(2)\)},
Canad.~J.~Math.~31 (1979), 726--785. Dihedral subcase reduction.
\item[Deligne, P.]
\emph{Formes modulaires et repr\'esentations \(\ell\)-adiques},
S\'em.~Bourbaki~355 (1971). Construction of the \(\ell\)-adic Galois
representation attached to a holomorphic newform \(F\) of weight
\(k\geq 2\); prerequisite for automorphy-lifting.
\item[Deligne, P. and Serre, J.-P.]
\emph{Formes modulaires de poids~\(1\)}, Ann.~Sci.~\'ENS.~7 (1974), 507--530.
\item[Ribet, K. A.]
\emph{Galois representations attached to eigenforms with nebentypus},
Glasgow Math.~J.~27 (1977). CM form classification.
\item[Selberg, A.]
\emph{Old and new conjectures and results about a class of Dirichlet
series}, Collected Papers II, Springer 1989, 47--63. Selberg class and
orthogonality.
\item[Luo, W., Rudnick, Z., and Sarnak, P.]
\emph{On the generalised Ramanujan conjecture for \(\mathrm{GL}(n)\)},
Proc.~Symp.~Pure Math.~66 (1999), 301--310.
\item[Iwaniec, H.]
\emph{Spectral Methods of Automorphic Forms} (second edition), AMS
Graduate Studies 53, 2002.
\item[M\oe glin, C. and Waldspurger, J.-L.]
\emph{Le spectre r\'esiduel de \(\mathrm{GL}(n)\)},
Ann.~Sci.~\'ENS~22 (1989), 605--674.
\end{description}

\paragraph{Disjoint-source discipline.}
The unconditional derivation catalogue (Kim--Shahidi 2002, Kim 2003,
Kim--Sarnak 2003, JPSS 1983, Jacquet--Shalika 1981) is disjoint from
the Route~D verification catalogue (Arthur 2013, Langlands--Shelstad
1987, Ng\^o 2010, Waldspurger 2008, Chaudouard--Laumon 2012,
Ramakrishnan 2000, Cogdell--Piatetski-Shapiro 1999, Labesse--Langlands
1979). The two catalogues share only JPSS 1983 and Jacquet--Shalika
1981, used structurally (defining \(L\)-functions and establishing
non-vanishing on \(\mathrm{Re}(s)=1\)), not to establish bounds.
Route~C depends on the Maass symmetric-power conjecture (open); as a
conditional, no disjointness check is required for its verification
source.
